%% ===================================================================
\section{Preliminaries and Threat Model}
\label{sec:background}
%% ===================================================================

Let $G = (V, E, A)$ be an undirected graph with $N = |V|$ nodes, adjacency $A \in \R^{N \times N}$, features $X \in \R^{N \times d_{\mathrm{in}}}$, and symmetric normalized adjacency $\Ahat = D^{-1/2}(A + I)D^{-1/2}$. Throughout, $\phi$ is the activation, $\sigma_i(\cdot)$ the $i$-th singular value, $\mathrm{vec}(\cdot)$ column-major vectorization. The sensitivity matrices ($S$, $S_c$, $S_K$, $S_v$) and governing scalars ($\kappa{=}\norm{J_z}_2$, $\rho(J_z)$, $\eta$, $\ecrit$, $r_v$) are defined where they first appear.

\textbf{IGNN and IFT.} \AEGIS analyses models that settle to an equilibrium, because that fixed point is what exposes structural sensitivity in closed form. A DEQ-GNN~\cite{bai2019deep,bai2020multiscale} iterates a weight-tied $F_\theta$ to a fixed point; the IGNN instantiation~\cite{gu2020implicit} uses
\begin{equation}
F(Z, A) = \phi(\Ahat Z W^\top + X_{\mathrm{proj}}),
\label{eq:ignn_operator}
\end{equation}
with $W \in \R^{d \times d}$ spectral-normalized~\cite{miyato2018spectral} so $\kappa \leq \norm{\Ahat}_2 \norm{W}_2 < 1$, ReLU $\phi$, and a learned projection $X_\mathrm{proj}$. Training uses implicit differentiation through the fixed point~\cite{gould2021deep,fung2022jfb} with optional Jacobian regularization~\cite{bai2021stabilizing}, and a linear head $f(\zstar) = W_\mathrm{cls}\zstar + b$ produces class predictions. At equilibrium, $G(\zstar,A){=}\zstar{-}F(\zstar,A){=}0$ defines $\zstar$ as an implicit function of $A$, with IFT sensitivity
\begin{equation}
\frac{\partial \zstar}{\partial p} = (I - J_z)^{-1} \frac{\partial F}{\partial p}\bigg|_{\zstar}.
\label{eq:ift}
\end{equation}
Equation~\eqref{eq:ift} is the lever \AEGIS pulls: it gives the closed-form first-order change in $\zstar$ under any parameter $p$, with the resolvent $(I-J_z)^{-1}$ amplifying a perturbation by at most $1/(1-\kappa)$ via its convergent Neumann series. The pseudospectral index $\eta{=}\norm{(I-J_z)^{-1}}_2(1-\rho) \geq 1$~\cite{trefethen2005spectra} measures the gap to the spectral-radius bound ($\eta{=}1$ for normal $J_z$).

\textbf{Threat model.} A white-box attacker with full access to the trained IGNN perturbs the normalized adjacency
\begin{equation}
\Ahat' = \Ahat + \delta \Ahat, \quad \norm{\delta \Ahat}_F \leq \varepsilon,
\label{eq:threat}
\end{equation}
with $\delta\Ahat$ symmetric ($\delta\Ahat{=}\delta\Ahat^\top$, preserving the undirected graph), edge-only ($[\delta\Ahat]_{ij}{=}0$ for $(i,j)\notin E$, so no new edges), and continuous in $\R$. \AEGIS targets edge-deletion / weight-perturbation; insertion attacks (Nettack-class~\cite{zugner2018adversarial,zugner2019adversarial} and node-injection~\cite{sun2020nipa}) require enlarging the basis to $\bar E \supseteq E$, a natural follow-up that preserves the $S_c$ construction. Discrete deletions connect to the continuous model through \cref{prop:transfer}. Unlike input-feature perturbation~\cite{el2021implicit,revay2020lipschitz}, structural perturbation changes the operator itself: $\zstar$ depends on $A$ through both $J_z = \mathrm{diag}(\phi') (\Ahat \otimes W)$ and $J_A = \partial F / \partial \mathrm{vec}(A)$, motivating the $S_c$ construction of \cref{sec:framework}.
